\chapter{Introducción}

\paragraph{Relación con la guía teórica.} Este documento es el \emph{compañero
numérico} de la guía \emph{Cálculo de Malliavin: de los procesos estocásticos y la
integral de Itô al estado del arte} (\codigo{docs/pdf/Calculo\_de\_Malliavin\_corregido.pdf}),
que contiene las demostraciones completas. Aquí se adopta su notación y se cita su
numeración con la forma \guia{Teorema 8.9}. Cada resultado se acompaña de su
implementación y de una figura reproducible. La correspondencia capítulo a capítulo está en
\codigo{docs/mapa\_pdf.md}.

El cálculo de Malliavin es un cálculo diferencial en dimensión infinita sobre el
espacio de Wiener, introducido por P.~Malliavin \cite{malliavin1978} para dar una
prueba probabilística del teorema de hipoelipticidad de Hörmander. Su herramienta
central es una \emph{fórmula de integración por partes} que permite trasladar
derivadas desde una función hacia un \emph{peso} aleatorio:
\[
  \E\big[\varphi'(F)\,G\big] = \E\big[\varphi(F)\,H(F,G)\big].
\]
En finanzas esta identidad tiene tres usos principales, que recorre este documento:
\begin{enumerate}
  \item \textbf{Cobertura} (cap.~\ref{cap:clark-ocone}): la fórmula de
        Clark--Ocone identifica explícitamente la cartera replicante
        \cite{ocone1984,oconekaratzas1991}.
  \item \textbf{Sensibilidades} (cap.~\ref{cap:griegas}): las griegas se calculan
        por Monte Carlo sin derivar el pago, lo que es crucial para pagos
        discontinuos \cite{fournie1999,fournie2001,benhamou2003}.
  \item \textbf{Densidades y volatilidad implícita} (cap.~\ref{cap:densidades}), además de esperanzas condicionales y opciones bermudas \cite{bally2005}:
        existencia y estimación de densidades \cite{nualart2006} y desarrollos
        asintóticos de la volatilidad implícita \cite{alos2007,alosgarcialorite2021}.
\end{enumerate}

Cada resultado se acompaña de una implementación en la librería Python
\codigo{malliavin} (directorio \codigo{src/malliavin}) y de un script de
visualización en \codigo{examples/}; las figuras se incluyen desde
\codigo{figures/}. El capítulo~\ref{cap:implementacion} describe el mapa
fórmula~$\to$~función.

\paragraph{Notación.} $W=(W_t)_{t\in[0,T]}$ es un movimiento browniano estándar en
$(\Omega,\F,\Prob)$ con filtración natural $(\F_t)$; $\Hil=L^2([0,T])$ con producto
$\ip{h}{g}=\int_0^T h(t)g(t)\dd t$; $W(h)=\int_0^T h\dd W$ es la integral de Wiener.
